\section{Diskusjon}

\subsection{Tolkning av grunnfrekvensen}
Ved alle tre turtallene var det grunnfrekvensen som ga størst utslag. En dominerende 1X-topp passer med ubalanse eller lignende \parencite{vibrasjon2}. Amplituden økte merkbart fra 0{,}346\,mm/s ved 3998\,rpm til 2{,}239\,mm/s ved 6001\,rpm. Dette passer med at ubalansekrefter øker når turtallet øker \parencite{vibrasjon2}.

\subsection{Vurdering av lagerfrekvensene}
At lagerfrekvensene er lave, betyr ikke at alle mulige feil kan utelukkes. Vi finner likevel ikke grunnlag for å påvise lokale skader i ytterring, innerring eller kule \parencite{vibrasjon3}. De lave amplitudene ved lagerfrekvensene tilsier tilstandsgrad 1 med hensyn til lokal lagerskade \parencite{tbv2}.

\subsection{Vurdering etter ISO 10816}
Målingene kan sammenlignes indikativt med ISO 10816-3 for små maskiner, der grenseverdiene er 0{,}71\,mm/s, 1{,}80\,mm/s og 4{,}50\,mm/s \parencite{iso10816, vibrasjon1}. Ved 2498 og 3998\,rpm ligger amplitudene ved grunnfrekvensen i sone A, mens 6001\,rpm havner i sone C. Det er økningen ved grunnfrekvensen som drar opp verdien, ikke lagerfrekvensene. Sammenligningen må tolkes med forsiktighet, siden ISO-vurdering normalt bygger på samlet RMS-verdi og ikke enkeltamplituder ved bestemte frekvenser.

\subsection{Feilkilder}
Usikkerheten ligger først og fremst i målekjeden. Sensorplassering og innfesting kan påvirke hvor mye av vibrasjonen som faktisk fanges opp, særlig dersom signalet dempes mellom lageret og målepunktet \parencite{tbv2}. I tillegg kan bakgrunnsstøy fra riggen gjøre små topper vanskelige å skille fra støygulvet. Frekvensoppløsningen i FFT-analysen kan også ha påvirket hvor presist frekvenstoppene ble avlest.

Målingene er dessuten et øyeblikksbilde. I praksis bygger prediktivt vedlikehold på trendanalyse over tid, der endringer i vibrasjonsnivå over flere målinger gir et mer pålitelig bilde \parencite{vibrasjon2, tbv2}. Siden forsøket ble gjort på en kontrollert laboratorierigg, er typiske feilkilder som feilsmøring, forurensning og feilmontering mindre aktuelle her \parencite{vibrasjon3}.

% TODO: Legg inn frekvensspekter-figurer fra SKF Observer som støtter diskusjonen
